% Laget automatisk av elevark/lag_elevark.py fra opplegg/repetisjon-felt.html. Endringer her blir overskrevet.
\documentclass[11pt]{article}
\usepackage{elevark}
\begin{document}
\elevarkhode{Fysikk 2 · Elektriske og magnetiske felt}{Repetisjon: elektriske og magnetiske felt}{https://lektorbo.no/opplegg/repetisjon-felt.html}{Opplegget på nett}

Elektriske og magnetiske felt ligner på hverandre, men de virker på ladninger på ulike måter. Her er en oversikt over likheter og forskjeller, de viktigste formlene og blandede oppgaver.

\section*{Elektrisk og magnetisk felt side om side}

\begin{center}\small
\begin{tabularx}{\linewidth}{YYY}
\toprule
\textbf{} & \textbf{Elektrisk felt $E$} & \textbf{Magnetisk felt $B$} \\
\midrule
Lages av & ladninger & ladninger i bevegelse (strøm) og magneter \\
Enhet & N/C = V/m & $T$ (tesla) \\
Kraft på ladning & $F = qE$, også når ladningen står i ro & $F = qvB \sin  θ$, bare når ladningen beveger seg \\
Retning på kraften & langs feltet (positiv ladning) eller mot feltet (negativ) & vinkelrett på både $v$ og $B$ (høyrehåndsregelen) \\
Arbeid på ladningen & ja, endrer farten: $W = qU$ & nei, endrer bare retningen \\
Bane når partikkelen kommer inn vinkelrett på et homogent felt & parabel, som et vannrett kast & sirkel med $r = mv/(qB)$ \\
Felt rundt en punktladning eller rett leder & $E = kq/r^{2}$, avtar som $1/r^{2}$ & $B = μ_{0}I/(2πr)$, avtar som $1/r$ \\
Feltlinjer & starter på positive og slutter på negative ladninger & lukkede kurver uten start og slutt \\
\bottomrule
\end{tabularx}
\end{center}

\section*{Formler}

\begin{itemize}
\item Coulombs lov: $F = k |q_{1}q_{2}| / r^{2}$, med $k = 8{,}99 \cdot 10^{9}\,\mathrm{N\,m^{2}/C^{2}}$.
\item Elektrisk felt: $E = F/q$. Rundt en punktladning: $E = k|q|/r^{2}$. Mellom to parallelle plater: $E = U/d$.
\item Akselerasjon gjennom spenningen $U$ fra ro: $qU = \tfrac{1}{2}mv^{2}$.
\item Magnetisk kraft på ladning: $F = qvB \sin  θ$. Sirkelbane: $r = mv/(qB)$ og $T = 2πm/(qB)$.
\item Fartsvelger: $qE = qvB$ gir $v = E/B$.
\item Magnetisk kraft på leder: $F = IlB \sin  θ$.
\item Magnetfelt rundt rett leder: $B = μ_{0}I/(2πr)$, med $μ_{0} = 4π \cdot 10^{-7}\,\mathrm{T\,m/A}$. Kraft per meter mellom to parallelle ledere: $F/l = μ_{0}I_{1}I_{2}/(2πd)$.
\end{itemize}

Konstanter: $e = 1{,}602 \cdot 10^{-19}\,\mathrm{C}$, $m_{\mathrm{e}} = 9{,}109 \cdot 10^{-31}\,\mathrm{kg}$, $m_{\mathrm{p}} = 1{,}673 \cdot 10^{-27}\,\mathrm{kg}$, $u = 1{,}6605 \cdot 10^{-27}\,\mathrm{kg}$ og $g = 9{,}81\,\mathrm{m/s^{2}}$.

\section*{Oppgaver}

Oppgavene blander elektriske og magnetiske felt. Simuleringene i fotnotene kan du bruke til å sjekke svarene.

\begin{oppgaveliste}
\oppgave \fase{forklar}{Begrep}Er påstandene sanne eller usanne? Begrunn svaret.

a) Et magnetfelt kan øke farten til et proton.\newline b) Et elektrisk felt kan øke farten til et proton.\newline c) Et nøytron blir ikke avbøyd i et magnetfelt.\newline d) Magnetiske feltlinjer starter på nordpoler og slutter på sydpoler, også inne i magneten.\newline e) Et elektron som ligger i ro i et magnetfelt, får ingen kraft.
\svar
\oppgave \fase{test}{Regn}En ladning på +3,0~nC ligger i $x = 0$, og en ladning på −1,0~nC ligger i $x = 10\,\mathrm{cm}$. a) Regn ut kraften mellom ladningene. b) Regn ut feltstyrken midt mellom ladningene, med retning. c) Finn punktet på x-aksen der feltet er null. Sjekk med simuleringen.\footnote{\url{lektorbo.no/sim/elektrisk-felt.html}}
\svar
\oppgave \fase{test}{Regn}Et elektron starter i ro ved den negative av to parallelle plater med 2,0~cm mellom og 500~V spenning. a) Finn feltstyrken og kraften på elektronet. b) Hvor stor fart har elektronet når det treffer den positive plata? c) Hvor lang tid bruker det?
\svar
\oppgave \fase{test}{Regn}En kule med massen 0,50~g og ladningen +20~nC henger i en snor mellom to plater med 6,0~cm mellom. Snora danner 10° med loddlinjen. Finn den elektriske kraften, feltstyrken og spenningen mellom platene. Sjekk med simuleringen.\footnote{\url{lektorbo.no/sim/ladd-kule.html}}
\svar
\oppgave \fase{test}{Regn}Et proton med farten $2{,}0 \cdot 10^{6}\,\mathrm{m/s}$ kommer vinkelrett inn i et magnetfelt på 0,20~T. Finn kraften på protonet, radien i banen og omløpstiden. Sjekk med simuleringen.\footnote{\url{lektorbo.no/sim/magnetfelt-bane.html}}
\svar
\oppgave \fase{test}{Regn}Neonioner med ladning +e går først gjennom en fartsvelger med $E = 4{,}0 \cdot 10^{4}\,\mathrm{V/m}$ og $B = 0{,}10\,\mathrm{T}$, og så inn i et magnetfelt på 0,50~T i et massespektrometer. a) Hvilken fart har ionene som kommer gjennom fartsvelgeren? b) Hvor langt fra hverandre treffer ${}^{20}\mathrm{Ne}$ (19,992~u) og ${}^{22}\mathrm{Ne}$ (21,991~u) detektoren etter en halvsirkel? Fartsvelger\footnote{\url{lektorbo.no/sim/fartsvelger.html}} og massespektrometer\footnote{\url{lektorbo.no/sim/magnetfelt-bane.html}}.
\svar
\oppgave \fase{test}{Regn}En 15~cm lang metallstav med massen 12~g ligger vannrett i et vannrett magnetfelt på 0,25~T som står vinkelrett på staven. Hvor stor strøm må gå i staven for at den magnetiske kraften skal kunne bære tyngden? Tegn en figur som viser retningen på strømmen, feltet og kraften. Se også strømhuska.\footnote{\url{lektorbo.no/sim/stromhuske.html}}
\svar
\oppgave \fase{test}{Regn}To lange, parallelle ledere står 8,0~cm fra hverandre. Den ene har 6,0~A, og den andre har 10~A i motsatt retning. Finn magnetfeltet midt mellom lederne, og kraften per meter mellom dem. Tiltrekker eller frastøter de hverandre? Sjekk med simuleringen.\footnote{\url{lektorbo.no/sim/magnetfelt-ledere.html}}
\svar
\oppgave \fase{utfordring}{Sammensatt}Et elektron akselereres fra ro gjennom 200~V og kommer så vinkelrett inn i et magnetfelt på 1,0~mT. a) Finn radien i banen. b) Vi legger til et elektrisk felt slik at elektronet går rett fram. Hvor sterkt må feltet være, og hvilken retning må det ha i forhold til kraften fra magnetfeltet? c) Hvorfor kan vi ikke bruke et magnetfelt til å akselerere elektronet?
\svar
\end{oppgaveliste}

\end{document}
